% 09_oq5_energy_resolution.tex — Section 9: OQ-5 Energy Resolution
% Phase 2 content migration — Task #178, Cycle 54 NVB-1, 2026-03-20
% Source: open_questions/oq5_energy_resolution.tex (rewrite 2026-03-20; authoritative)
% Law 16: quantitative claims carry source locator comments.
% Status: POPULATED (Phase 2) — key results summary.
%
% NOTE: Authoritative source is oq5_energy_resolution.tex.
%       The pre-rewrite version (48 eV NEP-based) is archived as
%       oq5_energy_resolution_pre_rewrite_2026-03-20.tex (incorrect model, retained for audit trail).

\section{OQ-5: Energy Resolution}
\label{sec:oq5}

\subsection{Detector Paradigm: Threshold Detector (Corrected)}

\textbf{The previous model was incorrect.} The pre-rewrite version applied a TES
linear calorimeter formula ($E_\mathrm{res} = \mathrm{NEP} \times \sqrt{\tau}$) and
obtained a phonon-noise floor of 48\,eV. This was wrong: D0 is a \emph{threshold
detector}, not a TES calorimeter.
% internal: oq5_energy_resolution.tex Sec.~1 (paradigm correction)

In a threshold detector, the superconducting film sits at $T_\mathrm{op} < T_c$.
When an energy deposition heats the pixel above $T_c$, the film transitions to normal
and the Meissner screening collapses (OQ-1 mechanism). The \emph{physically correct}
figure of merit is the \textbf{threshold energy}~\cite{Richards1994}:
% Richards1994 Sec. II.H p.17 (superconducting bolometers operate at the SC transition)
\begin{equation}
  \boxed{
  E_\mathrm{th} = \int_{T_\mathrm{op}}^{T_c} C_\mathrm{total}(T)\,dT
  \quad [\text{J}]
  }
  \label{eq:E_th}
\end{equation}
This is a \emph{thermodynamic} property of the pixel, not a noise-limited quantity.
An energy deposition $E_\mathrm{dep} < E_\mathrm{th}$ cannot trigger the transition
regardless of readout sensitivity.
% APPEC2021DirectDetection Sec. 4.6.1 p.28 (threshold model; cryogenic bolometer energy thresholds in sub-K systems)

\subsection{Pixel Heat Capacity Model}

The D0 pixel consists of:
\begin{enumerate}
  \item \textbf{NbN thin film} (Sommerfeld model): $C_\mathrm{NbN}(T) = \gamma_V V_\mathrm{NbN} T$,
    with $V_\mathrm{NbN} = 1.50 \times 10^{-18}$\,m$^3$ ($1.5\times 5.0\times 0.20\,\mu$m pixel).
    $\gamma_V \approx 313\,\text{J\,m}^{-3}\text{K}^{-2}$ (Drude estimate from
    Sidorova~\cite{Sidorova2020}).
    % Sidorova2020 Table III p.11 (N(0) = 4.0e28 eV^-1 m^-3 for disordered NbN films)
    % Richards1994 Appendix p.23 (gamma = pi^2/3 k_B^2 N(0), Sommerfeld model)
  \item \textbf{Ge absorber}: $C_\mathrm{Ge}(T) \propto V_\mathrm{Ge} T^3 / \Theta_D^3$
    (Debye low-T limit~\cite{Richards1994}), % Richards1994 Appendix p.22 (C = 234 N k_B (T/Theta)^3)
    with $V_\mathrm{Ge} = 7.50 \times 10^{-18}$\,m$^3$ ($1\,\mu$m thick) and
    $\Theta_D^\mathrm{Ge} = 374$\,K.
\end{enumerate}
The Ge absorber dominates total heat capacity ($\sim$80\%) because its volume is $5\times$
larger than the NbN film.
% internal: oq5_energy_resolution.tex Sec.~2 (Ge absorber dominates C_total)

\subsection{Threshold Energy Results}

Using $T_c = 15.7$\,K~\cite{Kamlapure2010}:
% Kamlapure2010 (NbN T_c ~ 15.87 K for epitaxial films; 15.7 K is D0 design value)

\begin{table}[ht]
\centering
\caption{Threshold energy $E_\mathrm{th}$ for the D0 pixel ($T_c = 15.7$\,K) vs.\
  operating temperature. NbN-only = no Ge absorber; Full pixel = NbN + $1\,\mu$m Ge.
  % internal: oq5\_energy\_resolution.tex Table 2 (E\_th vs. T\_op)
  Energies computed from Eqs.~(\ref{eq:E_th_NbN_sec9})--(\ref{eq:E_th_total_sec9}).}
\label{tab:E_th_table}
\small
\begin{tabular}{@{}p{2.5cm}p{3cm}p{3cm}p{4cm}@{}}
\toprule
$T_\mathrm{op}$ & NbN only & NbN + Ge & Note \\
\midrule
12.0\,K & 288\,keV & 1.56\,MeV & D0 operating point (current cryostat) \\
15.0\,K & 60\,keV & 383\,keV & $T_c - T_\mathrm{op} = 0.7$\,K \\
15.5\,K & 17.5\,keV & 114\,keV & $0.2$\,K below $T_c$ \\
15.68\,K & 1.76\,keV & 11.6\,keV & $20$\,mK below $T_c$ \\
\bottomrule
\end{tabular}
\end{table}

The analytic integrals are:
\begin{align}
  E_\mathrm{th}^\mathrm{NbN} &= \frac{\gamma_V V_\mathrm{NbN}}{2}
    \bigl(T_c^2 - T_\mathrm{op}^2\bigr), \label{eq:E_th_NbN_sec9} \\
  E_\mathrm{th}^\mathrm{Ge}  &= \frac{12\pi^4 n_\mathrm{Ge} V_\mathrm{Ge} k_B}
    {20\,\Theta_D^3} \bigl(T_c^4 - T_\mathrm{op}^4\bigr). \label{eq:E_th_Ge_sec9} \\
  E_\mathrm{th}^\mathrm{total} &= E_\mathrm{th}^\mathrm{NbN} + E_\mathrm{th}^\mathrm{Ge}. \label{eq:E_th_total_sec9}
\end{align}

\subsection{Key Finding: Operating Temperature is the Primary Lever}

At the D0 operating point ($T_\mathrm{op} = 12$\,K), $E_\mathrm{th} \approx 1.56$\,MeV
for the full NbN+Ge pixel. \textbf{This means: a 5.9\,keV Fe-55 X-ray will NOT trigger
a Meissner transition at $T_\mathrm{op} = 12$\,K.}

\textbf{The most effective design lever is operating temperature:} approaching $T_c$
dramatically reduces $E_\mathrm{th}$. At $T_\mathrm{op} = 15.5$\,K (0.2\,K below $T_c$),
$E_\mathrm{th} \approx 114$\,keV for the full pixel. The implication is that D0-Easy
(Meissner--Zeeman POC) is the correct first milestone: using a thermal sweep through $T_c$
rather than single-particle energy deposition.

\textbf{Also effective:} removing the Ge absorber reduces $E_\mathrm{th}$ by $5\times$
(NbN-only: 288\,keV vs.\ 1.56\,MeV full pixel at 12\,K).

\subsection{Model Limitations}

\needref{NbN thin-film Sommerfeld coefficient $\gamma_V$ from direct calorimetric measurement.
Estimate $\gamma_V \approx 200$--$600\,\text{J\,m}^{-3}\text{K}^{-2}$ depending on
film disorder and stoichiometry; Drude estimate used here: $\approx 313\,\text{J\,m}^{-3}\text{K}^{-2}$.}
The $\gamma_V$ uncertainty propagates directly into $E_\mathrm{th}$ (NbN-only). The Ge
contribution is insensitive to this uncertainty.
% internal: oq5_energy_resolution.tex Sec.~5 (limitations: gamma_V uncertainty)
